%% kapa: exam=1 round=31 year=2020 period=1 form=A subject=civil_law subject_ko=민법 seq=10 number=10
%% source: Q-Net 감정평가사 기출문제 (cst00309) artlSeq=5207219 file=extracted/1008/5207219/A형/A형/2020년도 제31회 감정평가사 1차 시험_1교시_A형.pdf
%% conversion: kordoc 4.24.51 -> tools/gen1.py (auto); review=proofread-sonnet; flags=none
\begin{question}{10}
대리에 관한 설명으로 옳지 않은 것은? (다툼이 있으면 판례에 따름)
\begin{choices}
  \item 계약체결의 권한을 수여받은 대리인은 체결한 계약을 처분할 권한이 있다.
  \item 본인이 이의제기 없이 무권대리행위를 장시간 방치한 것을 추인으로 볼 수는 없다.
  \item 매매계약의 체결과 이행에 관한 대리권을 가진 대리인은, 특별한 사정이 없으면 매수인의 대금지급기일을 연기할 수 있는 권한을 가진다.
  \item 본인이 사회통념상 대리권을 추단할 수 있는 직함이나 명칭 등의 사용을 승낙한 경우, 수권행위가 있는 것으로 볼 수 있다.
  \item 무권대리행위가 제3자의 위법행위로 야기된 경우에도, 본인이 추인하지 않으면 무권대리인은 계약을 이행하거나 손해를 배상하여야 한다.
\end{choices}
\end{question}
\answer{1}
